% part: intuitionistic-logic
% chapter: propositions-as-types
% section: introduction

\documentclass[../../../include/open-logic-section]{subfiles}

\begin{document}

\olfileid[hi]{int}{pty}{int}

\olsection{परिचय}

लैम्ब्डा कलन अभिकलन का सरल निदर्श है। यह चरों से बने पदों पर काम करता है। यदि
$N$ और $M$ पद हैं, तो $(NM)$ भी पद है (``$M$ पर लगाया गया~$N$'')। यदि $N$ पद
है, तो $\lambd[x][N]$ पद है। यदि $N$ में चर~$x$ है, तो $\lambd[x][N]$ में
~$x$ की संगत आवृत्तियाँ $\lambd[x]$ संकारक से आबद्ध होती हैं और इसलिए अब
मुक्त नहीं रहतीं। $(\lambd[x][N]M)$ रूप का पद, $\Subst{N}{M}{x}$ में
$\beta$-संकुचित होता है ($N$ में $x$ की हर मुक्त आवृत्ति को $M$ से बदलने का
परिणाम)। पद $N$,~$N'$ में $\beta$-रूपांतरित होता है यदि $N'$ किसी उपपद को
$\beta$-संकुचित करके $N$ से मिलता है; हम $N \redone N'$ लिखते हैं।
$\lambd$-कलन में अभिकलन कोई अनुक्रम $N \redone N_1 \redone N_2 \redone
\dots$ है। यदि यह अनुक्रम ऐसे पद $N_k$ पर समाप्त होता है जिसका कोई उपपद
$\beta$-संकुचित नहीं हो सकता, तो उसे \emph{प्रसामान्य} या~$N$ का प्रसामान्य
रूप कहते हैं। $\lambd$-कलन में अभिकलन को ऐसे अनुक्रमों के रूप में समझें।
प्रसामान्य रूप मिलने पर अभिकलन रुकता है और वही उसका परिणाम है। अभिकलनों का यह
सरल निदर्श ट्यूरिंग-पूर्ण निकलता है: प्रत्येक आंशिक पुनरावर्ती फलन की गणना
$\lambd$-कलन के किसी पद द्वारा की जा सकती है।

हास्केल करी और विलियम हावर्ड ने स्वतंत्र रूप से प्राकृतिक निगमन और
$\lambd$-कलन के बीच निकट समानता खोजी। सहजतः पद $\lambd[x][N]$ कोई फलन है:
$x$ उसका तर्क है और $N$ बताता है कि~$x$ दिए जाने पर मान कैसे निकाला जाए। तब
पद $(\lambd[x][N]M)$ फलन $\lambd[x][N]$ को तर्क~$M$ पर लगाना है, और
$\beta$-संकुचन बताता है कि उसका मूल्यांकन कैसे करें: $N$ में $x$ को~$M$ से
बदलें (और परिणामी पद का मूल्यांकन जारी रखें)। अब इन्हें नियत प्रांत और परिसर
वाले फलन मानें। यदि $\lambd[x][N]$ का प्रांत $A$ और परिसर~$B$ है, तो चर~$x$
को केवल~$A$ से तर्क लेने वाला और $N$ को~$B$ से मान उत्पन्न करने वाला समझना
चाहिए। फलतः $(\lambd[x][N]M)$ तभी अर्थपूर्ण होना चाहिए जब $\lambd[x][N]$
फलन $A \to B$ और $M \in A$ हो। यह सूचना $\lambd$-कलन में जोड़ने पर
\emph{प्रकारित} $\lambd$-कलन मिलता है। प्रकारित $\lambd$-कलन में प्रत्येक
व्यंजक $(\lambd[x][N]M)$ वैध नहीं, केवल वे जिनमें $\lambd[x][N]$ और $M$ के
``प्रकार'' मेल खाते हैं; अर्थात् $\lambd[x][N]$ का प्रकार $A \lif B$ और $M$
का प्रकार~$A$ है। $\lambd[x][N]$ का प्रकार $A \to B$ तभी है जब $x$ का प्रकार
$A$ और $N$ का प्रकार~$B$ हो। सामान्यतः हर व्यंजक $(M_1M_2)$ वैध नहीं है। केवल
वे पद $(M_1M_2)$ वैध हैं जहाँ $M_1$ का प्रकार $A \to B$ और $M_2$ का प्रकार
~$A$ है; तब $(M_1M_2)$ का प्रकार~$B$ है।

ये प्रकारण नियम अब स्पष्टतः प्राकृतिक निगमन की !!{derivation}s के अनुरूप हैं,
जहाँ प्रकारों को !!{formula}s और स्वयं !!{derivation}s को $\lambd$-पद निरूपित
करते हैं। उदाहरणतः $!A \lif !B$ की !!a{derivation} किसी पद
$\lambd[\typeof{x}{!A}][N]$ के अनुरूप है, जिसका फलन प्रकार $!A \lif !B$ है।
प्राकृतिक निगमन के अनुमान-नियम प्रकारित लैम्ब्डा कलन के प्रकारण नियमों के
अनुरूप हैं, उदाहरणतः
\begin{prooftree}
  \AxiomC{$\Discharge{!A}{x}$}
  \DeduceC{$!B$}
  \DischargeRule{\Intro\lif}{x}
  \UnaryInfC{$!A \lif !B$}
  \DisplayProof\bottomAlignProof
  \quad\text{के अनुरूप है}\quad
  \Axiom$x: !A \fCenter N: !B$
  \UnaryInf$ \fCenter \lambd[\typeof{x}{!A}][N] : !A \lif !B$
\end{prooftree}
जहाँ दाईं ओर के नियम का अर्थ है कि यदि $x$ का प्रकार~$!A$ और $N$ का
प्रकार~$!B$ है, तो $\lambd[\typeof{x}{!A}][N]$ का प्रकार~$!A \lif !B$ है।

$\Elim\lif$ नियम संयोजन पदों के प्रकारण नियम के अनुरूप है, अर्थात्
\[
  \AxiomC{$!A \lif !B$}
  \AxiomC{$!A$}
  \RightLabel{\Elim\lif}
  \BinaryInfC{$!B$}
  \DisplayProof
  \quad\text{के अनुरूप है}\quad
  \Axiom$\fCenter M_1: !A \lif !B$
  \Axiom$\fCenter M_2: !A$
  \BinaryInf$ \fCenter (M_1M_2) : !B$
  \DisplayProof
\]
यदि किसी \Intro\lif{} नियम के तुरंत बाद \Elim\lif{} नियम आता है, तो
!!{derivation} को सरल किया जा सकता है:
\begin{gather*}
  \AxiomC{$\Discharge{!A}{x}$}
  \DeduceC{$!B$}
  \DischargeRule{\Intro\lif}{x}
  \UnaryInfC{$!A \lif !B$}
  \AxiomC{}
  \DeduceC{$!A$}
  \RightLabel{\Elim\lif}
  \BinaryInfC{$!B$}
  \DisplayProof
  \quad\redone\quad
  \AxiomC{}
  \DeduceC{$!A$}
  \DeduceC{$!B$}
  \DisplayProof
\end{gather*}
जो लैम्ब्डा पदों के $\beta$-संकुचन के अनुरूप है
\[
(\lambd[\typeof{x}{!A}][N]M) \quad\redone\quad
\Subst{N}{M}{x}.
\]

$\land$ के नियमों और ``गुणनफल'' प्रकारों के बीच, तथा $\lor$ के नियमों और
``योग'' प्रकारों के बीच, ऐसी ही अनुरूपताएँ हैं।

सरलतः प्रकारित लैम्ब्डा कलन के पदों और प्राकृतिक निगमन की !!{derivation}s के
बीच इस अनुरूपता को ``करी--हावर्ड'', ``प्रोग्रामों के रूप में प्रमाण'', या
``प्रकारों के रूप में प्रतिज्ञप्तियाँ'' अनुरूपता कहते हैं। !!{formula}s
(प्रतिज्ञप्तियों) के प्रकारों से और प्रमाणों के पदों से अनुरूप होने के अतिरिक्त,
अनुरूपताओं को इस प्रकार सारांशित कर सकते हैं:
\begin{center}
  \begin{tabular}{c c c}
    तर्क & प्रोग्राम \\
    \hline
    प्रतिज्ञप्ति/!!{formula} & प्रकार \\
    प्रमाण & पद \\
    मान्यता & चर \\
    मुक्त की गई मान्यता & आबद्ध चर \\
    मुक्त न की गई मान्यता & मुक्त चर \\
    $!A \lif !B$ & फलन प्रकार \\
    $!A \land !B$ & गुणनफल प्रकार \\
    $!A \lor !B$ & योग प्रकार \\
    $\lfalse$ & रिक्त प्रकार \\
    \hline
  \end{tabular}
\end{center}

करी--हावर्ड अनुरूपता स्वचालित प्रमाण सहायकों और प्रोग्रामों के प्रकार-जाँचकों
की आधारशिलाओं में से एक है, क्योंकि किसी प्रतिज्ञप्ति का साक्षी प्रमाण जाँचना
(जैसा हमने ऊपर किया) यह जाँचने के तुल्य है कि किसी प्रोग्राम (पद) का घोषित
प्रकार है या नहीं।
\end{document}
